\chapter{Mathematical induction}
\chaptercontents{A & Weak induction (\S4.2) \\ B & Strong induction and well-ordering (\S4.3) \\ C & Review set 4}%
{\ess{4.2: 5, 11; 4.E: 10, 11, 16, 17, 18, 20, 25}\\ \ess{4.3: 8, 9, 14, 15; 4.E: 13, 19, 26, 27, 29, 30, 31, 32}\\ \rec{4.2: 8, 17, 20; 4.E: 12, 14, 15, 21, 24, 33; 4.3: 10, 18; 4.E: 28, 34, 35}\\ (\S4.1, Peano's axioms, is not in the planning.)}

\hsection{Weak induction}

\begin{result}[Principle of (weak) induction]
Let $p(n)$ be a statement about natural numbers $n$ and let $b\in\N$. Suppose
\begin{itemize}
\item \textbf{(base case)} $p(b)$ is true, and
\item \textbf{(induction step)} for all $n\ge b$: if $p(n)$ is true then $p(n+1)$ is true.
\end{itemize}
Then $p(n)$ is true for all $n\ge b$. (Most often $b=0$ or $b=1$.)
\end{result}
In the induction step the assumption ``$p(n)$'' is the \term{induction hypothesis}. You prove $p(n+1)$ \emph{from} it, for an arbitrary $n\ge b$; you do not get to assume $p(n+1)$.

\begin{strategy}[The induction template (use these words)]
\begin{enumerate}
\item ``For $n\in\N$ (with $n\ge b$) let $p(n)$ be the statement: \dots''. \quad $p(n)$ is a \emph{statement}, not a number.
\item ``\textbf{Base case} $n=b$: \dots'' \ Verify $p(b)$ by direct computation.
\item ``\textbf{Induction step.} Let $n\ge b$ and suppose $p(n)$ holds, \ie \dots\ (IH). We prove $p(n+1)$, \ie \dots''\\
   Start from the $n+1$ expression, isolate the $n$ expression inside it, apply the IH, finish with algebra.
\item ``By induction, $p(n)$ holds for all $n\ge b$.'' $\square$
\end{enumerate}
\end{strategy}

\begin{key}[Notation that appears in the exercises]
$\sum_{k=1}^{n}a_k=a_1+\dots+a_n$ and $\prod_{k=1}^na_k=a_1\cdots a_n$; an empty sum is $0$, an empty product is $1$. Recursive definitions: $0!=1$, $(n+1)!=(n+1)\cdot n!$; a sequence given by $a_0$ and a rule for $a_{n+1}$ in terms of $a_n$. The key identity for the step is $\sum_{k=1}^{n+1}a_k=\big(\sum_{k=1}^na_k\big)+a_{n+1}$.
\end{key}

\begin{hexample}[a sum formula]
Prove that $\displaystyle\sum_{k=1}^{n}k=\frac{n(n+1)}2$ for all $n\in\N$.
\solution
For $n\in\N$ let $p(n)$ be the statement $\sum_{k=1}^nk=\frac{n(n+1)}2$.\\
\emph{Base case} $n=0$: the empty sum is $0$ and $\frac{0\cdot1}2=0$. So $p(0)$ holds.\\
\emph{Induction step.} Let $n\in\N$ and suppose $\sum_{k=1}^nk=\frac{n(n+1)}2$ (IH). Then
\[
\sum_{k=1}^{n+1}k=\Big(\sum_{k=1}^{n}k\Big)+(n+1)\overset{\text{IH}}{=}\frac{n(n+1)}2+(n+1)=\frac{(n+1)(n+2)}2=\frac{(n+1)\big((n+1)+1\big)}2,
\]
which is $p(n+1)$.\\
By induction, $p(n)$ holds for all $n\in\N$.\qed
\end{hexample}

\begin{hexample}[geometric sum]
Let $r\in\R$ with $r\ne1$. Prove that $\displaystyle\sum_{k=0}^{n}r^k=\frac{r^{n+1}-1}{r-1}$ for all $n\in\N$.
\solution
Base case $n=0$: LHS $=r^0=1$, RHS $=\frac{r-1}{r-1}=1$.\\
Induction step: let $n\in\N$ and assume the formula for $n$. Then
$\sum_{k=0}^{n+1}r^k=\frac{r^{n+1}-1}{r-1}+r^{n+1}=\frac{r^{n+1}-1+r^{n+2}-r^{n+1}}{r-1}=\frac{r^{n+2}-1}{r-1}$, the formula for $n+1$.\\
By induction the formula holds for all $n\in\N$.\qed
\end{hexample}

\begin{hexample}[divisibility]
Prove that $3\mid 4^n-1$ for all $n\in\N$.
\solution
Base case $n=0$: $4^0-1=0=3\cdot0$. \\
Induction step: let $n\in\N$ and suppose $3\mid4^n-1$; fix $k\in\Z$ with $4^n-1=3k$. Then
\[
4^{n+1}-1=4\cdot4^n-1=4(4^n-1)+3=4\cdot3k+3=3(4k+1),
\]
so $3\mid4^{n+1}-1$. By induction, $3\mid4^n-1$ for all $n\in\N$.\qed\\
\emph{Trick: write the $n+1$ quantity as (something)$\cdot$(the $n$ quantity) $+$ (something divisible).}
\end{hexample}

\begin{hexample}[an inequality with a larger base case]
Prove that $2^n>n^2$ for all $n\ge5$.
\solution
Base case $n=5$: $2^5=32>25=5^2$.\\
Induction step: let $n\ge5$ and suppose $2^n>n^2$. Then $2^{n+1}=2\cdot2^n>2n^2$, so it suffices to show $2n^2\ge(n+1)^2$, \ie $n^2-2n-1\ge0$, \ie $(n-1)^2\ge2$, which holds since $n-1\ge4$. Hence $2^{n+1}>(n+1)^2$.\\
By induction, $2^n>n^2$ for all $n\ge5$.\qed\\
\emph{Pattern for inequalities: bound the $n+1$ side by the IH, then prove a separate ``gap'' inequality for $n\ge b$.}
\end{hexample}

\begin{hexample}[closed form of a recursive sequence]
Let $a_0=1$ and $a_{n+1}=2a_n+1$ for $n\in\N$. Prove that $a_n=2^{n+1}-1$ for all $n\in\N$.
\solution
Base case: $a_0=1=2^1-1$. \ Induction step: let $n\in\N$ with $a_n=2^{n+1}-1$. Then $a_{n+1}=2a_n+1=2(2^{n+1}-1)+1=2^{n+2}-1$. By induction, $a_n=2^{n+1}-1$ for all $n$.\qed
\end{hexample}

\begin{key}[Binomial coefficients and the binomial theorem]
For $n,k\in\N$ with $k\le n$, $\binom nk=\dfrac{n!}{k!\,(n-k)!}$ is the number of $k$-element subsets of an $n$-element set ($\binom nk=0$ for $k>n$).\\
\textbf{Pascal's rule:} $\binom{n+1}{k+1}=\binom nk+\binom n{k+1}$. \quad \textbf{Binomial theorem:} for $x,y\in\R$ and $n\in\N$,
\[
(x+y)^n=\sum_{k=0}^n\binom nk x^k y^{n-k}.
\]
Both are proved by induction on $n$ (the theorem's step uses Pascal's rule). Consequences: $\sum_{k=0}^n\binom nk=2^n$ ($x=y=1$), $\sum_{k}(-1)^k\binom nk=0$ for $n\ge1$.
\end{key}

\begin{hexample}[Bernoulli's inequality]
Let $x\in\R$ with $x\ge-1$. Prove that $(1+x)^n\ge1+nx$ for all $n\in\N$.
\solution
Base case $n=0$: $(1+x)^0=1\ge1+0$. \ Induction step: let $n\in\N$ and suppose $(1+x)^n\ge1+nx$. Since $1+x\ge0$ we may multiply the IH by $1+x$ without flipping the inequality:
\[
(1+x)^{n+1}=(1+x)(1+x)^n\ge(1+x)(1+nx)=1+(n+1)x+nx^2\ge1+(n+1)x,
\]
as $nx^2\ge0$. By induction the inequality holds for all $n\in\N$.\qed
\end{hexample}

\begin{warn}
\begin{itemize}
\item No base case, or the wrong base case ($n=0$ when the claim is for $n\ge1$): the proof is invalid. The classic ``all horses have the same colour'' fallacy has a correct-looking step that fails exactly from $n=1$ to $n=2$.
\item Never start the step with the statement $p(n+1)$ and manipulate it to something true; start from one side of $p(n+1)$ and transform it, or show both sides equal a common expression.
\item State $p(n)$ explicitly and say where you use the IH. ``Assume true for $n$'' without saying what is true is not acceptable.
\item If the step uses $p(n-1)$ as well as $p(n)$, you are doing strong induction: see Section B for the extra base cases.
\end{itemize}
\end{warn}

\begin{planning}[4.2 and 4.E (weak induction)]
\ess{4.2: 5, 11; 4.E: 10, 11, 16, 17, 18, 20, 25}\quad \rec{4.2: 8, 17, 20; 4.E: 12, 14, 15, 21, 24, 33}\\[3pt]
Skills: sum/product formulas, divisibility, inequalities with a base case $b>0$, recursively defined sequences, statements about finite sets (\eg number of subsets), binomial coefficients and the binomial theorem. Every one of them is the template above with a different algebraic middle.
\end{planning}

\hsection{Strong induction and well-ordering}

\begin{result}[Principle of strong induction]
Let $p(n)$ be a statement about natural numbers and $b\in\N$. Suppose that for every $n\ge b$:
\[
\text{if } p(k) \text{ holds for all } k \text{ with } b\le k<n, \text{ then } p(n) \text{ holds.}
\]
Then $p(n)$ holds for all $n\ge b$.\\
In practice: prove the first few values $p(b),\dots,p(b+m)$ directly (\term{base cases}), then for $n>b+m$ prove $p(n)$ assuming $p(k)$ for \emph{all} $b\le k<n$ (or just for the specific smaller $k$ you need, such as $n-1$ and $n-2$).
\end{result}
Weak and strong induction are equivalent principles, but strong induction is the right tool whenever the ``reason'' for $p(n)$ involves a smaller value that is not $n-1$: factoring $n=ab$, the recurrence $f_{n}=f_{n-1}+f_{n-2}$, subtracting $4$ from $n$, dividing $n$ by $b$.

\begin{strategy}[Strong induction template]
\begin{enumerate}
\item ``Let $p(n)$ be the statement \dots''. Decide which smaller values the step will need, and hence how many base cases.
\item ``\textbf{Base cases.}'' Check each one directly.
\item ``\textbf{Induction step.} Let $n\ge\dots$ and suppose $p(k)$ holds for all $b\le k<n$ (IH). \dots\ Therefore $p(n)$.'' Say explicitly which $k<n$ you apply the IH to, and check that $k\ge b$.
\item ``By strong induction, $p(n)$ holds for all $n\ge b$.''
\end{enumerate}
\end{strategy}

\begin{hexample}[prime factorisation exists]
Prove that every natural number $n\ge2$ is a product of (one or more) primes.
\solution
Let $p(n)$: ``$n$ is a product of primes''. Let $n\ge2$ and suppose $p(k)$ for all $2\le k<n$ (IH).\\
\emph{Case 1:} $n$ is prime. Then $n$ is a product of one prime.\\
\emph{Case 2:} $n$ is not prime. Then $n=ab$ with $1<a,b<n$ (a positive divisor other than $1$ and $n$ exists). By the IH, $a$ and $b$ are products of primes, so $n=ab$ is too.\\
In both cases $p(n)$ holds, so by strong induction every $n\ge2$ is a product of primes.\qed\\
\emph{No separate base case was needed: for $n=2$ the IH is vacuous and Case 1 applies.}
\end{hexample}

\begin{hexample}[Fibonacci bound, two base cases]
Let $f_0=0$, $f_1=1$ and $f_{n+2}=f_{n+1}+f_n$. Prove that $f_n\le2^{\,n-1}$ for all $n\ge1$.
\solution
\emph{Base cases:} $f_1=1\le2^0$ and $f_2=1\le2^1$.\\
\emph{Step:} let $n\ge3$ and suppose $f_k\le2^{k-1}$ for all $1\le k<n$. Since $n-1\ge1$ and $n-2\ge1$ the IH applies to both:
\[
f_n=f_{n-1}+f_{n-2}\le2^{n-2}+2^{n-3}\le2^{n-2}+2^{n-2}=2^{n-1}.
\]
By strong induction, $f_n\le2^{n-1}$ for all $n\ge1$.\qed\\
\emph{Two base cases, because the step uses $n-2$ and is therefore only valid from $n=3$.}
\end{hexample}

\begin{hexample}[stamps: every $n\ge12$ is $4a+5b$]
Prove that every $n\in\N$ with $n\ge12$ can be written as $n=4a+5b$ with $a,b\in\N$.
\solution
\emph{Base cases:} $12=4\cdot3$, $13=4\cdot2+5$, $14=4+5\cdot2$, $15=5\cdot3$.\\
\emph{Step:} let $n\ge16$ and suppose the claim for all $12\le k<n$. Then $12\le n-4<n$, so $n-4=4a+5b$ for some $a,b\in\N$, and $n=4(a+1)+5b$.\\
By strong induction the claim holds for all $n\ge12$.\qed
\end{hexample}

\begin{hexample}[the division theorem, existence part]
Let $b\ge1$. Prove that for every $a\in\N$ there exist $q,r\in\N$ with $a=qb+r$ and $0\le r<b$.
\solution
By strong induction on $a$. Let $a\in\N$ and suppose the claim for all $a'<a$.\\
If $a<b$: take $q=0$, $r=a$. \ If $a\ge b$: then $0\le a-b<a$, so by the IH $a-b=q'b+r$ with $0\le r<b$; hence $a=(q'+1)b+r$.\qed\\
\emph{Uniqueness: if $qb+r=q'b+r'$ with $0\le r,r'<b$ then $b\mid r-r'$ and $|r-r'|<b$, so $r=r'$, then $q=q'$. Extend to negative $a$ via $-a$.}
\end{hexample}

\begin{hexample}[base-$b$ expansions exist]\label{ex:baseb}
Let $b\ge2$. Prove that every $n\ge1$ has a base-$b$ expansion.
\solution
Strong induction on $n$. Let $n\ge1$ and assume all $1\le k<n$ have expansions. Write $n=qb+d$ with $0\le d<b$ (division theorem). If $q=0$ then $n=d$ is a one-digit expansion. If $q\ge1$ then $1\le q<n$ (as $b\ge2$), so $q=(d_r\dots d_0)_b$ by the IH, and then $n=qb+d=(d_r\dots d_0\,d)_b$.\qed
\end{hexample}

\begin{key}[Well-ordering principle]
Every inhabited subset of $\N$ has a least element. This is equivalent to induction, and gives the \term{minimal counterexample} method: to prove $\forall n\in\N,\ p(n)$, suppose not; then $\{n\in\N\mid\neg p(n)\}$ is inhabited, so it has a least element $m$; derive a contradiction by producing a smaller counterexample (or by checking $m$ directly if $m$ is small).
\end{key}

\begin{hexample}[well-ordering]
Prove that there is no integer $n$ with $0<n<1$.
\solution
Suppose, for a contradiction, that $S=\{n\in\Z\mid 0<n<1\}$ is inhabited. Then $S\subseteq\N$, so by well-ordering $S$ has a least element $m$. From $0<m<1$ we get, multiplying by $m>0$, $0<m^2<m<1$. So $m^2\in S$ and $m^2<m$, contradicting the minimality of $m$.\qed
\end{hexample}

\begin{warn}
\begin{itemize}
\item Count the base cases: if the step uses $n-d$, you need base cases $b,b+1,\dots,b+d-1$ and the step starts at $n=b+d$. Always check that the smaller value you feed into the IH is still $\ge b$.
\item In the step, the IH is ``$p(k)$ for all $b\le k<n$'', and you prove $p(n)$ (not $p(n+1)$). Mixing the two conventions produces off-by-one errors.
\item ``$n$ is not prime, so $n=ab$'': say why $1<a,b<n$. Without that the IH does not apply.
\end{itemize}
\end{warn}

\begin{planning}[4.3 and 4.E (strong induction)]
\ess{4.3: 8, 9, 14, 15; 4.E: 13, 19, 26, 27, 29, 30, 31, 32}\quad \rec{4.3: 10, 18; 4.E: 28, 34, 35}\\[3pt]
Skills: recognising when weak induction is not enough; proofs with several base cases (Fibonacci-type recurrences, $n-d$ steps); existence of factorisations/expansions; proofs via the well-ordering principle and minimal counterexamples.
\end{planning}

\hsection{Review set 4}
\begin{multicols}{2}\small
\begin{itemize}
\item Weak: $p(b)$ and $\forall n\ge b,\ p(n)\Rightarrow p(n+1)$.
\item Strong: $\forall n\ge b,\ (\forall k,\ b\le k<n\Rightarrow p(k))\Rightarrow p(n)$; base cases as many as the step needs.
\item Template words: ``Let $p(n)$ be'', ``Base case'', ``Induction step'', ``(IH)'', ``By induction''.
\item Sums: $\sum_{1}^{n+1}=\sum_1^n+a_{n+1}$. Divisibility: $x_{n+1}=c\cdot x_n+(\text{multiple})$. Inequalities: IH then a gap inequality.
\item Binomial theorem $(x+y)^n=\sum\binom nkx^ky^{n-k}$; Pascal's rule.
\item Well-ordering: inhabited $S\subseteq\N$ has a least element; minimal counterexample.
\end{itemize}
\end{multicols}
